\section{Personas}
%
\subsection{Primäre Personas}
%
\begin{figure}[H]
    \centering
    \begin{minipage}[t]{0.48\textwidth}
        \centering
        \includegraphics[width=\linewidth]{personas_primary_1.png}
        \vspace{0.4em}

        \small (a) Persona: CS Student
    \end{minipage}
    \hfill
    \begin{minipage}[t]{0.48\textwidth}
        \centering
        \includegraphics[width=\linewidth]{personas_primary_2.png}
        \vspace{0.4em}

        \small (b) Persona: Data-Blog Poster
    \end{minipage}
    \caption{}
    \label{fig:primary_personas_1_2}
\end{figure}
\begin{figure}[H]
    \centering
    \includegraphics[scale=0.8]{personas_primary_3.png}
    \caption{Persona: Communications director}
    \label{fig:primary_pesona_3}
\end{figure}
%

\smallskip
Unsere primären Personas sind Nutzer mit variierenden technischen Fähigkeiten, aber mit einem gemeinsamen Interesse an Daten, Datenvisualisierung und Informationen. Unsere Dienstleistung wird privat, für Hobbys und beruflich verwendet. Die Nutzer haben einen hohen Anspruch an Genauigkeit und Korrektheit und wünschen sich, dass die Visualisierung so gut wie möglich ist. 
\subsection{Sekundäre Persona}
%
\begin{figure}[H]
    \centering
    \includegraphics[scale=0.8]{personas_secondary.png}
    \caption{Persona: Casual User}
    \label{fig:secondary_persona}
\end{figure}
%

\smallskip
Unsere sekundären Nutzer haben selbst vielleicht kein hohes Interesse an den Daten, nutzen aber unsere Dienstleistung, um trotzdem welche zu sammeln. Die App muss intuitiv und einfach genug sein, damit jeder Nutzer schnell zu den gewünschten Informationen / Visualisierungen kommt.
%
\subsection{Negative Persona}
\begin{figure}[H]
    \centering
    \includegraphics[scale=0.8]{personas_problem.png}
    \caption{Persona: problematischer User}
    \label{fig:problem_persona}
\end{figure}
%

\smallskip
Unsere zugespitzte Darstellung eines problematischen Nutzers. Daten können natürlich auch für negative Zwecke verwendet werden. Es gilt, persönliche Daten zu anonymisieren und zu schützen. Bezogene Daten müssen auf Richtigkeit überprüft werden, die Quellen müssen vertrauenswürdig und zuverlässig sein.
